\section{Vector spaces}\label{sec:vector_spaces}

This section is a natural continuation of \fullref{sec:modules}, however the choice of topics is more suitable for \fullref{ch:linear_algebra}.

\begin{remark}\label{rem:vector}
  The term \enquote{vector} is among the most overloaded in mathematics. For its most popular uses, \textcite[rem. B.1.6]{Sally2013Analysis} writes
  \begin{displayquote}
    What is a vector? That's easy to answer. A vector is an element of a vector space. Lots of people describe vector as a quantity having magnitude and direction.
  \end{displayquote}

  In \cite{MacTutor:EarliestUsesOfMathWordsV}, both words \enquote{vector} and \enquote{scalar} are attributed to \textcite{Hamilton2000Quaternions}, who introduces them specifically for \hyperref[def:quaternion_algebra]{quaternions}:
  \begin{displayquote}
    The algebraically \textit{real} part may receive, according to the question in which it occurs, all values contained on the one \textit{scale} of progression of number from negative to positive infinity; we shall call it therefore the \textit{scalar part}, or simply the \textit{scalar} of the quaternion \textellipsis
    \begin{equation*}
      \vdots
    \end{equation*}

    On the other hand, the algebraically \textit{imaginary} part, being geometrically constructed by a straight line, or radius vector, which has, in general, for each determined quaternion, a determined length and determined direction in space, may be called the \textit{vector part}, or simply the \textit{vector} of the quaternion; \textellipsis
    \begin{equation*}
      \vdots
    \end{equation*}

    We may therefore say that \textit{a quaternion is in general the sum of its own scalar and vector parts} \textellipsis
  \end{displayquote}

  Thus, in Hamilton's terminology, the quaternion \( q = a + bi + cj + dk \) has scalar part \( a \) and vector part \( bi + cj + dk \) (more conventionally represented by the \hyperref[def:ordered_tuple]{ordered triple} \( (b, c, d) \)).

  The notion of vector has evolved in several directions; we will compare them in an ambient \hyperref[def:affine_space]{affine space} \( (A, \vect A, \tau) \) over the \hyperref[def:field]{field} \( \BbbK \).
  \begin{thmenum}
    \thmitem{rem:vector/radius_vector} As can be seen from Hamilton's quote above, the term \enquote{radius vector} chronologically precedes \enquote{vector}.

    We define radius vectors in \cref{def:radius_vector} as oriented segments with a predefined origin, i.e. ordered pairs \( (O, x) \) of points of \( A \) where \( O \) is fixed.

    This definition is stated for the \hyperref[def:euclidean_plane]{Euclidean plane} in
    \cite[221]{Hadamard2008LessonsInGeometryVol1},
    \cite[86]{Тыртышников2007ЛинейнаяАлгебра},
    \cite[124]{Узков1951ВекторныеПространства} and
    \cite[69]{Станилов1974АналитичнаГеометрия}.

    \thmitem{rem:vector/bound_vector} We mention in \cref{def:directed_segment} that directed segments are also called \enquote{bound vectors}.

    These generalize radius vectors by not fixing an implicit starting point. For affine spaces, this usage can be found in \cite[\S 2.1.13]{Berger1987GeometryI}.

    \thmitem{rem:vector/free_vector} Free vectors represent, in the words of \textcite[1]{FriedbergInselSpence2018LinearAlgebra}, \enquote{entities involving both magnitude and direction}.

    Free vectors can be formalized via bound vectors, i.e. directed segments. We can declare two directed segments to be equivalent if they are \hyperref[def:affine_parallelism]{parallel}, \hyperref[def:geometric_ray/unidirectional]{unidirectional} and have equal \hyperref[def:line_segment_length]{lengths}.

    This induces an \hyperref[def:equivalence_relation]{equivalence relation} on bound vectors. The equivalence classes of these vectors are called \enquote{free vectors} in \cite[87]{Тыртышников2007ЛинейнаяАлгебра} and \cite[13]{Станилов1974АналитичнаГеометрия}.

    In the aforementioned books, free vectors are defined with respect to an informal ambient \hyperref[def:euclidean_plane]{Euclidean plane}. The construction is meaningless in affine spaces, where free vectors are taken as part of the definition. In fact, when introducing affine spaces. \textcite[\S 2.1.3]{Berger1987GeometryI} defines the \enquote{free vector} \( \vect{xy} \) associated with the \enquote{bound vector} \( (x, y) \) as the unique element of \( \vect A \) such that \( \tau_{\vect{xy}}(x) = y \).

    The term \enquote{free vector} is also used for elements of \( \vect A \) in \cite[def. II.9.1]{Bourbaki1998Algebra1to3}.

    \thmitem{rem:vector/vector_space} Among the authors we cite in this monograph, the dominant notion of \enquote{vector} is \enquote{an element of a vector space}, and the dominant notion of \enquote{scalar} is a value from the underlying field that acts on vectors by \enquote{scaling} them.

    \Fullref{thm:vector_space_basis_existence} implies that every vector space has a \hyperref[def:hamel_basis]{Hamel basis}, so, modulo the choice of \hyperref[def:ordered_hamel_basis]{ordered basis}, every vector is an \hyperref[def:ordered_tuple]{ordered tuple} or, more generally, an \hyperref[def:indexed_family]{indexed family} over the base field.

    Vector spaces can be generalized to \hyperref[def:module]{modules} and \hyperref[def:semimodule]{semimodules}, but, as shown in \cref{ex:def:hamel_basis/module_without_basis}, a basis may not exist. It is thus unclear whether their elements can rightfully be called \enquote{vectors}. For example, \textcite[78]{Мальцев2009ОсновыЛинейнойАлгебры} uses \enquote{vector} for elements of any module, while \textcite[308]{Aluffi2009Algebra} mentions that the term is only sometimes used for \hyperref[def:free_semimodule]{free modules}.

    Many authors discuss both vector spaces and more general modules, but only use the word \enquote{vector} in relation to vector spaces. This is done, for example, in
    \cite[34]{Knapp2016BasicAlgebra},
    \cite[118]{Lang2002Algebra},
    \cite[308]{Aluffi2009Algebra},
    \cite[247]{Rotman2015AdvancedModernAlgebraPart1},
    \cite[def. II.1.2]{Bourbaki1998Algebra1to3} and
    \cite[13]{Кострикин2020АлгебраЧасть2}.

    For modules, the above books still mostly use \enquote{scalar} and related terms like \enquote{scalar multiplication}.

    In books dedicated to commutative algebra, such as \cite{Eisenbud1995CommutativeAlgebra} and \cite{КоцевСидеров2016КомутативнаАлгебра}, the \enquote{scalar} and \enquote{vector} terminology is not used.

    On the other hand, in books featuring only linear algebra, \enquote{vector} is actively used for any element of an abstract vector space. This is done, for example, in
    \cite[7]{FriedbergInselSpence2018LinearAlgebra},
    \cite[185]{Курош1968КурсВысшейАлгебры} and
    \cite[111]{Обрешков1962ВисшаАлгебра}.

    \thmitem{rem:vector/column} \enquote{Vector} may also refer to \enquote{column vector}, i.e. real or complex \hyperref[def:array/matrix]{matrix} with only one column. This is a particular special case of \enquote{element of a vector space} for \( \BbbR^{n \times 1} \) or \( \BbbC^{n \times 1} \).

    This usage can be found in books dedicated to concrete linear algebra, for example in
    \cite[def. 1.1]{Stewart1998MatrixAlgorithmsVol1} and \cite[37]{Тыртышников2007ЛинейнаяАлгебра}.

    The dual notion of \enquote{row vectors} then corresponds to \hyperref[def:dual_vector_space]{linear functionals}.

    \thmitem{rem:vector/list} Finally, \enquote{vector} is sometimes used as a synonym for \enquote{\hyperref[def:ordered_tuple]{ordered tuple}}.

    This usage can be found in literature for computer programming and related disciplines. For example, \cite[724]{AbelsonSussman2012eSICP} write
    \begin{displayquote}
      Abstractly, a vector is a compound data object whose individual elements can be accessed by means of an integer index in an amount of time that is independent of the index.
    \end{displayquote}

    For accessing individual elements, \fullref{thm:cartesian_product_universal_property} provides us with \hyperref[def:cartesian_product_projection]{projection functions}\fnote{Coordinate projections can be viewed as a special case of projection maps in Cartesian products --- see \cref{def:basis_decomposition}.}. In programming, the ability to access an element at an arbitrary index is called \enquote{random access}. In the context of \hyperref[def:turing_machine]{Turing machines}, \textcite{AroraBarak2009ComputationalComplexity} write
    \begin{displayquote}
      \textit{Random access} denotes the ability to access the \( i \)th symbol of the memory within a single step, without having to move a head all the way to the \( i \)th location. The name \enquote{random access} is somewhat unfortunate since this concept involves no notion of randomness --- perhaps \enquote{indexed access} would have been better. However, \enquote{random access} is widely used, and so we follow this convention this book.
    \end{displayquote}

    For C++, the \Verb|std::vector| class template determines not only access, but also several mutation operations. The following overview is given in \cite[\S 24.3.11]{ISO:14882:2024}:
    \begin{displayquote}
      A \Verb|vector| is a sequence container that supports (amortized) constant time insert and erase operations at the end; insert and erase in the middle take linear time. Storage management is handled automatically, though hints can be given to improve efficiency
    \end{displayquote}

    In some restricted domain-specific languages, there are dedicated types for only short vectors. For example, the GLSL documentation in \cite{GLSL:DataType} lists \Verb|vec2|, \Verb|vec3| and \Verb|vec4|.

    In type theory, \enquote{vector} is used as a particular formalization of ordered tuples in \cite[170]{UnivalentFoundationsProgram2013HoTT} and \cite[\S 6.4.7]{Mimram2020ProgramEqualsProof}.
  \end{thmenum}
\end{remark}

\paragraph{Vector spaces}

\begin{definition}\label{def:vector_space}\mcite[78]{Мальцев2009ОсновыЛинейнойАлгебры}
  A \term[ru=линейно пространство (\cite[111]{Обрешков1962ВисшаАлгебра}), ru=векторное/линейное пространство, en=vector/linear space (\cite[7]{FriedbergInselSpence2018LinearAlgebra})]{vector space} is a \hyperref[def:module]{left module}\fnote{It is irrelevant if we use left or right modules because the notions are equivalent in commutative rings.} over a \hyperref[def:field]{field}.

  As per \fullref{thm:vector_space_basis_existence}, in classical mathematics, all vector spaces have a \hyperref[def:hamel_basis]{Hamel basis}. Furthermore, \cref{thm:vector_space_dimension} shows that all bases are \hyperref[def:equinumerosity]{equinumerous}, leading to a well-defined notion of dimension --- see \cref{def:vector_space_dimension}.

  There is little nuance to the definition itself, except for minor terminological differences like the ubiquity of the vector/scalar terminology from \cref{rem:vector}. We compare the basic metamathematical notions for completeness:
  \begin{thmenum}
    \thmitem{def:vector_space/theory}\mimprovised For vector spaces over \( \BbbK \), we can simply reuse the \hyperref[def:fol_equational_theory]{equational theory} for left \( \BbbK \)-modules from \cref{def:module/theory}.

    \thmitem{def:vector_space/morphism}\mcite[138]{Мальцев2009ОсновыЛинейнойАлгебры} For the \hyperref[def:module/morphism]{module homomorphisms}, we prefer the linear function terminology from \cref{def:linear_function}.

    \thmitem{def:vector_space/category}\mcite[example III.5.5]{Aluffi2009Algebra} We denote the category of vector spaces over \( \BbbK \) by \( \ucat[\BbbK]{Vect} \) rather than \( \ucat[\BbbK]{Mod} \).

    \medskip

    \thmitem{def:vector_space/subobject}\mcite[115]{Мальцев2009ОсновыЛинейнойАлгебры} We call the \hyperref[def:module/subobject]{submodules} of a vector space \term[bg=подпространство (\cite[101]{Обрешков1962ВисшаАлгебра}), ru=линейное подпространство, en=vector subspace (\cite[17]{FriedbergInselSpence2018LinearAlgebra})]{subspaces}.

    \thmitem{def:vector_space/generated}\mcite[30]{FriedbergInselSpence2018LinearAlgebra} For the \hyperref[def:module/generated]{generated submodules}, we prefer the linear span definition and terminology from \cref{def:linear_span} (characterizing the generated submodule via \hyperref[def:linear_combination]{linear combinations}).

    \thmitem{def:vector_space/trivial}\mcite[116]{Мальцев2009ОсновыЛинейнойАлгебры} We call the \hyperref[def:module/trivial]{zero (sub)module} the \term[ru=нулевое подпространство, en=zero subspace (\cite[17]{FriedbergInselSpence2018LinearAlgebra})]{zero (sub)space}.

    \thmitem{def:vector_space/kernel}\mcite[67]{FriedbergInselSpence2018LinearAlgebra} For the \hyperref[def:module/kernel]{kernel} of a linear function, another established term is \term[ru=нуль-пространство (матрицы) (\cite[112]{Тыртышников2007ЛинейнаяАлгебра}), en=null space (\cite[67]{FriedbergInselSpence2018LinearAlgebra})]{null space}.

    \thmitem{def:vector_space/quotient}\mcite[\S 26.4]{Тыртышников2007ЛинейнаяАлгебра} Unsurprisingly, we call the \hyperref[def:module/quotient]{quotient modules} of a vector space \term[bg=фактор-пространство (\cite[308]{ГеновМиховскиМоллов1991Алгебра}), ru=фактор-пространство, en=quotient space (\cite[exerc. 1.3.31(d)]{FriedbergInselSpence2018LinearAlgebra})]{quotient spaces}.
  \end{thmenum}
\end{definition}
\begin{comments}
  \item Because a field is a commutative ring, we may just as well define a vector space as a right module.
\end{comments}

\begin{proposition}\label{thm:def:vector_space}
  A \hyperref[def:vector_space]{vector space} \( V \) over \( \BbbK \) has the following basic properties:
  \begin{thmenum}
    \thmitem{thm:def:vector_space/span_independent} For a \hyperref[def:linear_dependence]{linearly independent} set \( A \), if \( x \in V \setminus \lin A \), then \( A \cup \set{ x } \) is also a linearly independent set.

    A related result is \cref{thm:def:linear_dependence/span_dependent}.

    \thmitem{thm:def:vector_space/maximal_independent} Every maximally linearly independent subset of \( V \) is a basis.

    The converse holds more generally --- see \cref{thm:def:hamel_basis/maximal_independent}.

    \thmitem{thm:def:vector_space/minimal_spanning} Every minimal spanning subset of \( V \) is a basis.

    The converse holds more generally --- see \cref{thm:def:hamel_basis/minimal_spanning}.

    \thmitem{thm:def:vector_space/expansion} Let \( U \) be a subspace of \( V \). Suppose that \( u_1, \ldots, u_m \) is a basis of \( U \) and \( v_1, \ldots, v_n \) --- of \( V \). Then \( m \leq n \) and there exist indices \( i_{m+1}, \ldots, i_n \) such that the following is a basis of \( V \):
    \begin{equation*}
      u_1, \ldots, u_m, v_{i_{m+1}}, \ldots, v_{i_n}
    \end{equation*}
  \end{thmenum}
\end{proposition}
\begin{proof}
  \SubProofOf{thm:def:vector_space/span_independent} Suppose that \( A \) is linearly independent and fix a vector \( x \) from \( V \setminus \lin A \).

  Let
  \begin{equation*}
    t_0 x + \sum_{k=1}^n t_k a_k = 0_M.
  \end{equation*}

  If \( t_0 \) is nonzero, we can divide by it to obtain
  \begin{equation*}
    x = -\sum_{k=1}^n \frac {t_k} {t_0} a_k,
  \end{equation*}
  which implies that \( x \) is in the span of \( A \). But this contradicts our choice of \( x \).

  It follows that \( t_0 = 0_\BbbK \), hence \( \sum_{k=1}^n t_k a_k = 0_M \) belongs to the span of \( A \) and thus
  \begin{equation*}
    t_1 = \cdots = t_n = 0_\BbbK.
  \end{equation*}

  It follows that any linear combination of \( A \cup \set{ x } \) that sums to zero must be trivial.

  \SubProofOf{thm:def:vector_space/maximal_independent} Let \( A \) be a maximally linearly independent set of \( V \).

  Suppose that it is not a spanning set and fix a value \( x \) from \( V \setminus \lin A \). By \cref{thm:def:vector_space/span_independent}, the set \( A \cup \set{ x } \) is linearly independent, contradicting the maximality of \( A \).

  Therefore, \( A \) is a spanning set for \( V \).

  \SubProofOf{thm:def:vector_space/minimal_spanning} Let \( A \) be a minimal spanning set of \( V \). Suppose that it is linearly dependent. Then there exist distinct vectors \( x_1, \ldots, x_n \) in \( A \) and scalars \( t_1, \ldots, t_n \), at least one of which is nonzero, such that
  \begin{equation*}
    \sum_{k=1}^n t_k x_k = 0_M.
  \end{equation*}

  Let \( k_0 \) be the smallest index such that \( t_{k_0} \neq 0_R \). Then
  \begin{equation*}
    x_{k_0} = -\sum_{k \neq k_0} \frac {t_k} {t_{k_0}} a_k.
  \end{equation*}

  Then \( \lin A = \lin A \setminus \set{ x_{k_0} } \), contradicting the minimality of \( A \).

  \SubProofOf{thm:def:vector_space/expansion} Let \( U \) be a subspace of \( V \), let \( u_1, \ldots, u_m \) be a basis of \( U \) and \( v_1, \ldots, v_n \) --- of \( V \).

  We will recursively build new bases of \( V \) until one contains \( u_1, \ldots, u_m \). More precisely, at step \( k = 0, 1, \ldots m\), we should have a basis
  \begin{equation}\label{eq:thm:def:vector_space/expansion/proof/old_basis}
    u_{m-k+1}, \ldots, u_m, v_{i_{k,k+1}}, \ldots, v_{i_{k,n}},
  \end{equation}
  where the base step starts from \( v_1, \ldots, v_n \) by setting \( i_{0,j} \coloneqq j \). As we will see, this procedure should terminate \emph{after} step \( k = m - 1 \) (i.e. \eqref{eq:thm:def:vector_space/expansion/proof/old_basis} will hold for \( k = m \)).

  Suppose we have already constructed a basis for step \( k \). Take a decomposition of \( u_{m-k} \):
  \begin{equation}\label{eq:thm:def:vector_space/expansion/proof/umk}
    u_{m-k} \coloneqq \sum_{j=1}^k t_j u_{m-k+j} + \sum_{j=k+1}^n t_j v_{i_{k,j}}.
  \end{equation}

  At least one of \( t_{k+1}, \ldots, t_n \) must be nonzero --- otherwise \( u_{m-k} \) will be a linear combination of \( u_{m-k+1}, \ldots, u_m \), which would contradict our assumption that they are linearly independent. Let \( j_0 > k \) be the smallest such index. Then
  \begin{equation*}
    v_{i_{k,j_0}} = \frac 1 {t_{j_0}} \parens*{ u_{m-k} - \sum_{j=1}^k t_j u_{m-k+j} - \sum_{j \neq j_0} t_j v_{i_{k,j}} }.
  \end{equation*}

  Therefore, \( v_{i_{k,j_0}} \) belongs to the linear span of
  \begin{equation}\label{eq:thm:def:vector_space/expansion/proof/new_basis}
    u_{m-k}, u_{m-k+1}, \ldots, u_m, v_{i_{k,k+1}}, \ldots, v_{i_{k,j_0-1}}, v_{i_{k,j_0+1}}, \ldots, v_{i_{k,n}}.
  \end{equation}

  We have already shown that \eqref{eq:thm:def:vector_space/expansion/proof/new_basis} spans \( V \). To conclude that it is a basis, we must show that the vectors are linearly independent. Indeed, take some linear combination
  \begin{equation*}
    0_M = s_0 u_{m-k} + \sum_{j=1}^k s_j u_{m-k+j} + \sum_{j \neq j_0} s_j v_{i_{k,j}}
  \end{equation*}
  of the vectors from \eqref{eq:thm:def:vector_space/expansion/proof/new_basis} that sums to \( 0_M \) and substitute \( u_{m-k} \) for its decomposition \eqref{eq:thm:def:vector_space/expansion/proof/umk}. This gives us a linear combination of basis vectors from \eqref{eq:thm:def:vector_space/expansion/proof/old_basis} that sums to \( 0_M \), which can only be trivial. It follows that \( s_0 = s_1 = \cdots = s_n = 0_\BbbK \).

  Therefore, \eqref{eq:thm:def:vector_space/expansion/proof/new_basis} is a linearly independent spanning set for \( V \), i.e. a basis. We renormalize the indices as
  \begin{equation*}
    i_{k+1,j} \coloneqq \begin{cases}
      i_{k+1,j},   &j < j_0, \\
      i_{k+1,j+1}, &j \geq j_0,
    \end{cases}
  \end{equation*}

  The proposition is proven after we reach step \( k = m \) (we cannot proceed further anyway). If we reach step \( k = n \), the vectors \( v_1, \ldots, v_n \) will be replaced by \( u_{m-n+1}, \ldots, u_m \), and the latter will be a basis of \( V \). In particular, \( u_{m-n+1}, \ldots, u_m \) will span the entirety of \( V \), making \( u_1, \ldots, u_{m-n} \) redundant. Unless \( n = m \), this will contradict the minimality of Hamel bases shown in \cref{thm:def:hamel_basis/minimal_spanning}.

  In particular, it follows that \( m \) cannot exceed \( n \).
\end{proof}

\paragraph{Ranks and dimensions}

\begin{theorem}[Vector space basis existence]\label{thm:vector_space_basis_existence}
  Every \hyperref[def:vector_space]{vector space} has a \hyperref[def:hamel_basis]{basis}.
\end{theorem}
\begin{comments}
  \item Within \hyperref[def:zfc]{\logic{ZF}}, this theorem is equivalent to the \hyperref[def:zfc/choice]{axiom of choice} --- see \cref{thm:axiom_of_choice_equivalences/vector_space_bases}.
\end{comments}
\begin{proof}
  \ImplicationSubProof[thm:zorns_lemma]{Zorn's lemma}[thm:vector_space_basis_existence]{vector space existence} Let \( V \) be a vector space over \( \BbbK \). Let \( \mscrB \) be the family of all linearly independent \hyperref[def:linear_combination]{subsets} of \( V \). This family is nonempty because, by \cref{thm:def:linear_dependence/empty}, the empty set is by itself linearly independent.

  The union of any chain \( \mscrB' \subseteq \mscrB \) can then contain only linearly independent elements since otherwise we would have that some set in \( \mscrB' \) is not linearly independent. Thus, \fullref{thm:zorns_lemma} shows the existence of a maximal linearly independent set \( E \). By \cref{thm:def:vector_space/maximal_independent}, \( E \) is a basis.

  \ImplicationSubProof[thm:vector_space_basis_existence]{vector space existence}[def:zfc/choice]{the axiom of choice} Shown in \cite{Blass1984BasesImplyAOC}.
\end{proof}

\begin{proposition}\label{thm:vector_space_dimension}
  All \hyperref[def:hamel_basis]{bases} of a \hyperref[def:vector_space]{vector space} are \hyperref[def:equinumerosity]{equinumerous}.
\end{proposition}
\begin{proof}
  Let \( A \) and \( B \) be bases of \( V \).

  \SubProof{Proof for finite bases} Suppose that \( a_1, \ldots, a_n \) are the vectors of \( A \).

  \SubProof*{Proof in case \( B \) is infinite} Aiming at a contradiction, suppose that \( B \) is infinite.

  Every vector \( a_i \) of \( A \) can be \hyperref[def:basis_decomposition]{decomposed} along \( B \) as
  \begin{equation*}
    a_i = \sum_{j=1}^{m_i} t_{i,j} b_{i,j}.
  \end{equation*}
  for appropriate scalars from \( \BbbK \) and vectors from \( E \). Since this can be done for every \( a_i \) in \( A \), we conclude that
  \begin{equation*}
    V = \lin\set{ a_1, \ldots, a_n } = \lin\set{ b_{1,1}, \ldots, b_{1,m_1}, \ldots, b_{n,1}, \ldots, b_{n,m_n} }.
  \end{equation*}

  Hence, a finite subset of \( E \) spans \( V \), contradicting \cref{thm:def:hamel_basis/minimal_spanning}. Therefore, \( E \) must be a finite set.

  \SubProof*{Proof in case \( E \) is finite} Let \( b_1, \ldots, b_m \) be the vectors of \( B \). By applying \cref{thm:def:vector_space/expansion} twice, we conclude that \( n \leq m \) and \( m \leq n \), hence \( n = m \).

  \SubProof{Proof for infinite bases} Suppose that both \( A \) and \( B \) are infinite. For every vector \( x \), define the set
  \begin{equation*}
    B_x \coloneqq \set{ b \in B \given \pi_b(x) \neq 0_\BbbK },
  \end{equation*}
  where \( \pi_b \) are the \hyperref[def:basis_decomposition]{coordinate projections}. Fix a \hyperref[def:choice_function]{choice function} on \( \seq{ B_a }_{a \in A} \) that chooses a basis vector \( b \) for every \( a \).

  Let \( \seq{ a_\beta }_{\beta < \alpha} \) be a transfinite enumeration of \( A \). We use \fullref{thm:bounded_transfinite_recursion} on \( \beta \) to build the family of \hyperref[def:set_valued_map/partial]{partial functions}
  \begin{equation*}
    \iota_\beta: A \to B \times \omega,
  \end{equation*}
  where \( \iota_\beta \) is the union of \( \iota_\gamma \) for \( \gamma < \beta \) extended with
  \begin{equation*}
    \iota_\beta(a_\beta) \coloneqq \parens[\big]{ c(a), \sharp_\beta(c(a)) }.
  \end{equation*}

  Here the helper \( \sharp_\beta(b) \) counts the \enquote{previous} occurrences of \( b \):
  \begin{equation*}
    \sharp_\beta(b) \coloneqq \card(\set{ \gamma < \beta \given \qexists* {n \in \omega} \iota_\gamma(a_\gamma) = (b, n) }).
  \end{equation*}

  This ensures that \( \iota_\beta \) is injective. Thus, \( \iota_\alpha: A \to B \times \omega \) is a total single-valued function that is by construction injective.

  Therefore,
  \begin{equation*}
    \card(A)
    \leq
    \card(B \times \omega)
    \reloset {\ref{thm:cardinality_product_rule}} =
    \card(B) \times \aleph_0
    \reloset {\ref{thm:simplified_cardinal_arithmetic/infinite}} \leq
    \card(B).
  \end{equation*}

  The converse inclusion follows analogously.
\end{proof}

\begin{definition}\label{def:vector_space_dimension}\mimprovised
  We define the \term[bg=размерност (\cite[216]{Станилов1974АналитичнаГеометрия}), ru=размерность (\cite[sec. 3.7]{Тыртышников2007ЛинейнаяАлгебра})]{dimension} \( \dim V \) of the \hyperref[def:vector_space]{vector space} \( V \) as the \hyperref[thm:cardinality_existence]{cardinality} of any of its bases.
\end{definition}
\begin{defproof}
  \Fullref{thm:vector_space_basis_existence} shows that every vector space has a basis. \Cref{thm:commutative_module_rank} shows that all bases are equinumerous, hence the rank is well-defined.
\end{defproof}

\begin{proposition}\label{thm:commutative_module_rank}
  All \hyperref[def:hamel_basis]{Hamel bases} of a \hyperref[def:module]{module} over a \emph{nonzero} \hyperref[def:ring/commutative]{commutative ring} are \hyperref[def:equinumerosity]{equinumerous}.
\end{proposition}
\begin{proof}
  Suppose that \( A \) and \( B \) are bases of the \( R \)-module \( M \).

  Let \( I \) be a \hyperref[def:semiring_ideal/maximal]{maximal ideal} of \( R \). By \cref{thm:quotient_by_maximal_ideal}, \( \BbbK = R / I \) is a field.

  Since \( A \) and \( B \) are both bases of \( M \), there exists some isomorphism \( \Phi: R^{\oplus A} \to R^{\oplus B} \) with components \( \Phi_b: R^{\oplus A} \to R \), which we can extend to
  \begin{equation*}
    \begin{aligned}
      &\Psi: \BbbK^{\oplus A} \to \BbbK^{\oplus B}, \\
      &\Psi\parens*{ \seq{ t_a + I }_{a \in A} } \coloneqq \seq*{ \Phi_b(\seq{ t_a }_{a \in A}) + I }_{b \in B}.
    \end{aligned}
  \end{equation*}

  This is clearly an isomorphism if it is well-defined. To show that it is well-defined, suppose \( t_a + I = t_a' + I \) for every \( a \) in \( A \), i.e. \( t_a - t_a' \) is in \( I \). Then
  \begin{equation*}
    \Phi_b(\seq{ t_a }_{a \in A}) - \Phi_b(\seq{ t_a' }_{a \in A})
    =
    \Phi_b(\seq{ t_a - t_a' }_{a \in A})
    =
    \sum_{a \in A} (\underbrace{ t_a - t_a' }_{\in I}) \Phi_b(a)
    \in
    I.
  \end{equation*}

  Thus,
  \begin{equation*}
    \Phi_b(\seq{ t_a }_{a \in A}) + I = \Phi_b(\seq{ t_a' }_{a \in A}) + I.
  \end{equation*}

  Therefore, \( \Psi \) is a well-defined linear isomorphism. Applying \cref{thm:vector_space_dimension}, we conclude that \( A \) and \( B \) are equinumerous.
\end{proof}

\begin{corollary}\label{thm:finitely_generated_module_basis}
  If a \hyperref[def:module/generated]{finitely generated} \( R \)-module over a nonzero commutative ring \( R \) has a \hyperref[def:hamel_basis]{basis}, this basis is finite.
\end{corollary}
\begin{proof}
  Follows from \cref{thm:commutative_module_rank}.
\end{proof}

\begin{definition}\label{def:module_rank}\mimprovised
  For a \hyperref[def:ring/commutative]{commutative ring} \( R \), we define the \term[bg=rank (\cite[32]{КоцевСидеров2016КомутативнаАлгебра}), ru=rank (\cite[47]{Шафаревич2019ОсновныеПонятияАлгебры}), en=rank (\cite[171]{Jacobson1985BasicAlgebraI})]{rank} of a \hyperref[def:free_semimodule]{free \( R \)-module} as the \hyperref[thm:cardinality_existence]{cardinality} of any of its bases.
\end{definition}
\begin{comments}
  \item Unlike vector space dimensions, module ranks may not exist in more general cases --- see \cref{ex:def:hamel_basis/module_without_basis}.
  \item We generalize the definition by \textcite[171]{Jacobson1985BasicAlgebraI} to infinite ranks.
\end{comments}
\begin{defproof}
  It follows from \cref{thm:commutative_module_rank} that the choice of basis is immaterial.
\end{defproof}

\begin{example}\label{ex:field_submodules}
  For a nonzero \hyperref[def:ring/commutative]{commutative ring} \( R \), every ideal of \( R \) is a submodule of the \hyperref[def:module_rank]{rank-one} \hyperref[def:module]{module} \( R \). For example, both \( \BbbZ \) and \( 2\BbbZ \) are \( \BbbZ \)-modules of rank one.

  For a \hyperref[def:field]{field} \( \BbbK \), every ideal of \( \BbbK \) is a submodule of the \hyperref[thm:vector_space_dimension]{unidimensional} \hyperref[def:vector_space]{vector space} \( \BbbK \). The only ideals of \( \BbbK \) are the zero ideal, whose dimension is zero, and the field itself, whose dimension is one.
\end{example}

\begin{proposition}\label{thm:modules_with_same_rank_are_isomorphic}
  Two \( R \)-modules having the same \hyperref[def:module_rank]{rank} are \hyperref[def:linear_function]{linearly isomorphic}.
\end{proposition}
\begin{proof}
  Suppose that \( A \) is a \hyperref[def:hamel_basis]{basis} of \( M \) and \( B \) is a basis of \( N \). Also suppose that both \( M \) and \( N \) have the same rank.

  Then \( A \) is equinumerous with \( B \) and there exists a bijective function \( f: A \to B \). We can define the map \( \varphi: A \to R^{\oplus B} \) by sending \( a \) to the embedding of \( f(a) \) in \( R^{\oplus B} \). We can then use \fullref{thm:free_semimodule_universal_property} to extend \( \varphi \) to a linear isomorphism \( \Phi: R^{\oplus A} \to R^{\oplus E} \).

  We now have isomorphisms between \( M \) and \( R^{\oplus A} \), between \( R^{\oplus A} \) and \( R^{\oplus B} \) and between \( R^{\oplus B} \) and \( N \). Composing them gives us an isomorphism between \( M \) and \( N \).
\end{proof}

\begin{proposition}\label{thm:rank_of_direct_sum}
  Assuming that the \( R \)-modules \( M_1, \ldots, M_n \) have bases, the \hyperref[def:module_rank]{rank} of their direct sum
  \begin{equation*}
    M_1 \oplus \cdots \oplus M_n
  \end{equation*}
  is the \hyperref[def:cardinal_arithmetic/addition]{sum of cardinals}
  \begin{equation*}
    \rank M_1 + \cdots + \rank M_n.
  \end{equation*}
\end{proposition}
\begin{proof}
  Follows from \cref{thm:basis_of_direct_sum}.
\end{proof}

\paragraph{Rank-nullity theorem}

\begin{lemma}\label{thm:vector_space_dimension_monotonicity}
  Let \( U \) be a subspace of the \hyperref[thm:vector_space_dimension]{finite-dimensional} \hyperref[def:vector_space]{vector space} \( V \). Then
  \begin{equation*}
    \dim U \leq \dim V.
  \end{equation*}
\end{lemma}
\begin{proof}
  Let \( A \) be a basis of \( U \). By \cref{thm:def:vector_space/expansion}, there exists a basis \( B \) of \( V \) such that \( A \subseteq B \). Then clearly \( \dim U \leq \dim V \).
\end{proof}

\begin{proposition}\label{thm:rank_nullity_via_subspaces}
  Let \( U \) be a subspace of the \hyperref[thm:vector_space_dimension]{finite-dimensional} \hyperref[def:vector_space]{vector space} \( V \). Then
  \begin{equation*}
    V \cong U \oplus (V / U).
  \end{equation*}
\end{proposition}
\begin{proof}
  By \cref{thm:vector_space_dimension_monotonicity}, \( \dim U \leq \dim V \). Let \( a_1, \ldots, a_m \) be an ordered basis for \( U \). By \cref{thm:def:vector_space/expansion}, there exist vectors \( a_{m+1}, \ldots, a_n \) such that \( a_1, \ldots, a_n \) is a basis of \( V \). We will show that \( a_{m+1} + U, \ldots, a_n + U \) is a basis of \( V / U \), thus demonstrating the isomorphism.

  First, we must show that it is a spanning set. Fix a vector \( x + U \) from \( V / U \). We know that \( x \) can be decomposed as
  \begin{equation*}
    x = \sum_{k=1}^n \pi_k(x) \cdot a_k.
  \end{equation*}

  Since \( a_1, \ldots, a_m \) belong to \( U \),
  \begin{equation*}
    x + U = \sum_{k={m+1}}^n \pi_k(x) \cdot a_k + U.
  \end{equation*}

  Hence, \( x + U \) can be represented as a linear combination of the vectors \( a_{m+1} + U, \ldots, a_n + U \).

  To show uniqueness, suppose that
  \begin{equation*}
    x + U = \sum_{k={m+1}}^n t_k (a_k + U) = \sum_{k={m+1}}^n r_k (a_k + U).
  \end{equation*}

  Then
  \begin{equation*}
    U = \sum_{k={m+1}}^n (t_k - r_k) (a_k + U).
  \end{equation*}

  Since neither of \( a_{m+1}, \ldots, a_n \) are in \( U \), and since they are linearly independent, it follows that \( t_k = r_k \) for \( k = m + 1, \ldots, n \).

  Therefore, \( a_{m+1} + U, \ldots, a_n + U \) is a basis of \( V / U \).
\end{proof}

\begin{definition}\label{def:rank_and_nullity}\mcite[70]{FriedbergInselSpence2018LinearAlgebra}
  For a \hyperref[def:linear_function]{linear map} \( \varphi: U \to V \) between \hyperref[def:vector_space]{vector spaces}, we call the \hyperref[def:vector_space_dimension]{dimension} of the \hyperref[def:vector_space/kernel]{kernel} of \( \varphi \) the \term[ru=дефект (\cite[250]{Тыртышников2007ЛинейнаяАлгебра})]{nullity} of \( \varphi \) and the dimension of the image - the \term[ru=ранг (\cite[250]{Тыртышников2007ЛинейнаяАлгебра})]{rank} of \( \varphi \).
\end{definition}

\begin{theorem}[Rank-nullity theorem]\label{thm:rank_nullity_theorem}
  For every \hyperref[def:linear_function]{linear map} \( \varphi: U \to V \) between \hyperref[thm:vector_space_dimension]{finite-dimensional} \hyperref[def:vector_space]{vector spaces}, we have
  \begin{equation}\label{eq:thm:rank_nullity_theorem}
    U \cong \ker \varphi \oplus \img \varphi.
  \end{equation}

  In particular,
  \begin{equation}\label{eq:thm:rank_nullity_theorem/dim}
    \dim U = \dim \ker \varphi + \dim \img \varphi.
  \end{equation}
\end{theorem}
\begin{proof}
  By \fullref{thm:homomorphism_theorem_for_modules},
  \begin{equation*}
    \img \varphi \cong U / \ker \varphi.
  \end{equation*}

  By \cref{thm:rank_nullity_via_subspaces},
  \begin{equation*}
    U \cong \ker \varphi \oplus (U / \ker \varphi) \cong \ker \varphi \oplus \img \varphi.
  \end{equation*}

  The equality \eqref{thm:rank_nullity_theorem} then follows from \cref{thm:rank_of_direct_sum}.
\end{proof}
